2. Data and Event Definition
2.1 PJM Load Data
Hourly electricity demand data were obtained from PJM Data Miner 2 (Hourly Load: Metered), the publicly accessible load reporting portal maintained by PJM Interconnection. All records use RTO-level aggregate load, representing the metered demand across the full PJM transmission system. The analysis period spans January 1, 2010 through December 31, 2014, yielding 43,824 hourly observations prior to feature construction. Following the creation of lagged and rolling-window features with a maximum lookback of 168 hours, the first 168 observations are dropped to avoid partially undefined feature rows, leaving 43,656 complete observations available for modeling.
The dataset is partitioned into a training set covering January 1, 2010 through December 31, 2013 (34,896 observations) and a held-out test set covering the full calendar year 2014 (8,760 observations). Model fitting and feature construction use only information available within the training period; no 2014 observations are included in training. The 2014 test year is used solely for evaluation.
Table 1 summarizes the load distribution across the training period, the 2014 test year, and the two event comparison windows analyzed in this study. The training set spans a wider load range (49,558–158,043 MW) than the 2014 test year (57,570–141,678 MW). The gap between the training maximum (158,043 MW) and the 2014 test maximum (141,678 MW) is noted as distributional context; the causes of this difference are not analyzed in this study. Within this context, the 2014 near-peak cold-event load of 140,510 MW, while near the 2014 annual peak, is not near the upper boundary of the training load distribution.
C03S — Revised Section 2.2 and Updated Notes

2.2 ERA5 Retrospective Weather Inputs (Revised)
Meteorological inputs are drawn from the ERA5 reanalysis dataset produced by the European Centre for Medium-Range Weather Forecasts (ECMWF) and distributed through the Copernicus Climate Data Store [CITATION: ERA5 reanalysis]. ERA5 provides hourly estimates of atmospheric variables on a global grid with approximately 31 km horizontal resolution, reconstructed through a numerical weather prediction model constrained by historical observations.
Three ERA5 near-surface variables are used: 2-metre air temperature (t2m), 10-metre zonal wind component (u10), and 10-metre meridional wind component (v10). Dewpoint temperature (d2m) was not included in the downloaded ERA5 files and is not used in the final modeling table.
Spatial aggregation uses the great_lakes_core subset, defined as a rectangular ERA5 grid box spanning 40.0°N–42.5°N and −88.0°W to −74.0°W at 0.25° resolution. The subset contains 627 grid cells, and each ERA5 variable is aggregated by a simple unweighted arithmetic mean over latitude and longitude. No population weighting is applied.
From these aggregated fields, derived features are constructed as follows. Temperature is converted from Kelvin to Fahrenheit. Wind speed is computed as sqrt(u10² + v10²) and converted from m/s to mph. Wind chill is computed using the NOAA/NWS wind chill formula when T ≤ 50°F and wind speed exceeds 3 mph; otherwise wind chill is set equal to temperature. Heating and cooling degree hours are computed as HDH = max(65 − T°F, 0) and CDH = max(T°F − 65, 0).
A critical interpretive constraint governs all results involving weather-informed models. ERA5 provides retrospective reanalysis values — meteorological conditions reconstructed after the fact using historical observations — rather than real-time operational numerical weather prediction (NWP) forecasts. All model results using ERA5 inputs therefore reflect an idealized retrospective weather-information setting and cannot be interpreted as achievable performance under operational deployment conditions, where forecast weather would be used in place of reanalysis weather. This distinction is preserved throughout the paper and in all comparisons involving the gradient boosted models.
NOAA ISD station data are not used in the modeling table and are not the primary source for any manuscript figures or conclusions.

2.3 Temporal Alignment and Feature-Ready Dataset
PJM load records are reported in Eastern Prevailing Time (EPT), which alternates between Eastern Standard Time (UTC−5) and Eastern Daylight Time (UTC−4) according to the U.S. daylight saving schedule. ERA5 weather data are natively UTC-timestamped. Temporal alignment is performed using UTC timestamps, with EPT used for reporting event times. This approach preserves the same-hour correspondence between meteorological conditions and observed load through daylight-saving transitions without introducing ambiguity at clock-change boundaries.
The transition from Eastern Daylight Time to Eastern Standard Time on November 2, 2014 produces a duplicate clock hour at 01:00 EPT. Both the 01:00 EDT and 01:00 EST records are preserved in the dataset with their respective UTC offsets, maintaining a continuous hourly series without data loss or gap insertion. Forecast evaluation at the November 2 transition hour uses the UTC-aligned weather record corresponding to each distinct EPT hour.
Lagged load features and rolling statistical features are constructed using the aligned hourly series. The maximum lookback horizon used in feature construction is 168 hours (one calendar week), which drives the 168-row trim at the start of the dataset described in Section 2.1. Feature construction details and the full feature set are described in Section 3.
2.4 Event Definition: January 6–8, 2014 Polar Vortex
The primary stress-test evaluation window is defined as the 72-hour period from 00:00 EPT January 6, 2014 through 23:00 EPT January 8, 2014, corresponding to the sustained cold air outbreak associated with the displacement of the stratospheric polar vortex over the eastern United States [CITATION: meteorological analysis of 2014 polar vortex]. The January 6–8, 2014 Polar Vortex produced a near-annual-peak RTO load of 140,510.2 MW at Jan 7 18:00 EPT, equal to 99.18% of the 2014 annual peak of 141,677.9 MW recorded on Jun 17 17:00 EPT.
This event window is defined as a near-annual-peak cold-weather stress event. It is not the 2014 annual peak, which occurred during the summer comparison window described in Section 2.5. No winter record or all-time peak designation is applied.
Table 1 characterizes the cold-event window in detail. Mean RTO load during the 72-hour window was 120,069 MW, substantially above the full-year test mean of 91,034 MW. The minimum observed load within the window was 82,638 MW, reflecting the sustained character of the cold outbreak rather than a brief isolated spike. Mean 2-meter temperature over the window was 3.4°F and mean HDH was 61.2 — the highest of any evaluation window summarized in Table 1. By comparison, the training period mean temperature was 52.3°F with a mean HDH of 12.7. This contrast between the meteorological conditions prevailing during training and those observed during the event window provides the basis for treating the polar vortex as an out-of-distribution generalization test for models fitted on 2010–2013 data.
The event is useful as a benchmarking case for several reasons beyond its meteorological severity. PJM RTO hourly load records for this period are publicly available through Data Miner 2 and ERA5 reanalysis inputs are publicly available through the Copernicus Climate Data Store, supporting a transparent retrospective benchmark against which future methods may be evaluated. The event was also extensively reviewed in post-event reports by NERC and PJM [CITATION: NERC Polar Vortex Review 2014; CITATION: PJM Polar Vortex Review 2014], providing independent documentation of the operational significance of the load levels observed.
Figure 1 summarizes the 2014 PJM load profile, the January 6–8 cold-event window, and the corresponding ERA5 retrospective weather conditions.
2.5 Summer Near-Peak Comparison Window
A summer comparison window is defined as the 72-hour period from 00:00 EPT June 16, 2014 through 23:00 EPT June 18, 2014. This window contains the 2014 annual peak of 141,678 MW, recorded at Jun 17 17:00 EPT, and serves as a structural counterpoint to the polar vortex window in the evaluation.
The two windows are selected to be comparable in terms of their proximity to the annual system peak — the polar vortex window reaches 99.18% of the annual peak — while representing opposing meteorological regimes. During the summer window, mean temperature was 73.3°F and mean HDH was 0.0, reflecting purely cooling-driven load rather than heating-driven load. Mean RTO load during the summer window was 126,845 MW with a minimum of 108,911 MW. Comparing model performance across these two near-peak windows — one winter-driven, one summer-driven — allows the evaluation to distinguish between general peak-load forecasting difficulty and the specific challenge of cold-weather extreme generalization. Results from this comparison are reported in Section 5.